Al estar el \textit{frontend} y el \textit{backend} desplegados como aplicaciones independientes, resultaba contraproducente recurrir al mecanismo tradicional de sesiones basado en \textit{cookies} del navegador. En su lugar, se optó por implementar un sistema de autenticación \textit{stateless} (sin estado en el servidor) basado en JSON Web Tokens (JWT)\footnote{\url{https://jwt.io/}} mediante la librería \texttt{djangorestframework-simplejwt}\footnote{\url{https://django-rest-framework-simplejwt.readthedocs.io/}}. Cuando un usuario inicia sesión, el servidor genera dos tokens, un \textit{Access Token} de corta duración (20 minutos) que el \textit{frontend} adjunta en cada petición para identificarse, y un \textit{Refresh Token} de larga duración (7 días) que permite renovar el primero sin obligar al usuario a volver a introducir sus credenciales.

Para gestionar los \textit{endpoints} de registro, inicio de sesión, cambio de contraseña y recuperación por correo electrónico se utiliza Djoser\footnote{\url{https://djoser.readthedocs.io/}}, una librería que proporciona todos estos flujos listos para usar sobre DRF, lo que reduce significativamente el tiempo de desarrollo y las posibles vulnerabilidades de seguridad.
